\section{Discussion}\label{sec:discussion}

\paragraph{Recommendations for clients}
Our results argue against both extremes clients could adopt. Approval should attach to what a tool can do, not to what it says. A client can compute a capability summary for each installed version, statically as we do or from a manifest the server must declare and the client verifies, and re-ask only when a high-risk capability is added. In our replay that rule prompts well under once per server per year. Clients should show the user what changed (``this tool can now run commands'') rather than a generic notice, and should treat changes to tool \emph{outputs} that instruct the model to act as part of the tool's behavior.

\paragraph{Recommendations for the protocol}
MCP annotations already let a server declare that a tool is read-only or non-destructive, but the declaration is self-reported and nothing ties it to the code or checks it across versions; in our data, hints stay unchanged through 96.2\% of high-risk capability gains in annotated tools (Section~\ref{sec:rq2}). The specification could make capability declarations mandatory for tools that touch the file system, processes, credentials, or the network, version them, and require clients to re-prompt when they grow. It should also state whether a server may instruct the agent to execute commands on its behalf; today some servers do and others explicitly forbid it.

\paragraph{Capabilities are not enough}
The email incident that motivates this work involved no new capability: the tool already sent email and simply sent it to one more address. Capability-based approval would not have caught it. Catching it requires describing where data may go (``only to recipients the user specified''), which our static analysis approximates only through new external hostnames. Data-destination declarations are a natural next step.

\paragraph{Limitations}
Our results are bounded in two ways. The first is scope. We analyze source code statically, so we establish that a path exists, not that it runs; capabilities hidden in dependencies, downloaded at run time, or reached through dynamic dispatch escape us, and our counts are lower bounds against an adversary who wants to hide. We study git histories, while users install packages from npm or PyPI that may differ. Sampling releases understates how often servers change. Our population is popular open-source servers in three languages, and eight of its repositories were deleted or made private during the study. The second is measurement. Labels come from a single annotator, so we cannot report agreement between humans, and because we label only flagged events, we measure precision but not recall.

\paragraph{Future work}
Comparing published registry packages with their source histories, adding data-flow destinations to the capability model, and measuring how users respond to capability-specific re-prompts are the obvious next steps.
